%% coverletter_template.tex -- starting point for every cover letter.
%% Copy, rename to  Lastname_Company_RoleShort_ReqID_CoverLetter.tex , fill,
%% compile with pdfLaTeX. Needs coverletter.cls beside it (or on TEXINPUTS)
%% and, for the script signature, signature.png beside it
%% (scripts/make_signature.py makes one). No signature.png = typed name.
%%
%% Angle-bracket text is placeholder. One page, always.
\documentclass{coverletter}
% \documentclass[typedname]{coverletter}   % also print the typed name under the signature

\name{<First Last>}
% \letterdate{<Month D, YYYY>}             % defaults to today's date
\contactline{<Street Address>}             % optional line
\contactline{<City, ST ZIP>}
\contactline{\href{mailto:<email>}{<email>}}
\contactline{<(000) 000-0000>}
\recipient{<Company> \\
           <City, ST> \\
           <Position Title (Req ID)>}

\begin{document}
\makeheader
\opening{Dear Hiring Manager:}

% Voice rules (full list in references/style_coverletter.md):
%   plain first person, light and personable, sounds like an engineer wrote it
%   no em dashes, no colons or semicolons inside the body
%   no flattery of the company, no mission-statement echo
%   hard facts live in the resume; pick one or two and tell them as a story
%   every fact must already be in the resume or the career profile

% P1: the role, who you are in one line, one human sentence about how you work.
<Paragraph 1.>

% P2: one concrete project from the resume that mirrors the posting's core work.
<Paragraph 2.>

% P3: name the honest gaps against the posting and bridge each one in a
% sentence or two. Drop this paragraph only if there are no real gaps, and
% use the space for a second relevant strength instead.
<Paragraph 3.>

% P4: people skills, logistics the posting asks about (citizenship, travel,
% clearance eligibility, relocation), and a plain invitation to talk.
<Paragraph 4.>

Thank you for your consideration. Please do not hesitate to contact me if you have any questions.

\closing{Sincerely,}
\end{document}
